% Pour toute question sur l'utilisation de cette classe, écrivez à jerome.soucy(at)ulaval.ca
% Nécessite la classe Evaluation.sty dans le même dossier ou dans le PATH de votre système
% Nécessite le fichier logo_ul.pdf dans le même dossier un dans le PATH de votre système
\def\Ponderation{35}
\def\SigleNumero{STT-2100}
\def\NomCours{Régression}
\def\DateEvaluation{19 décembre 2024}
\def\DateRemise{}
\def\HeureRemise{}
\def\HeureEvaluation{Durée: 1h50}
\def\Session{Été 2021}
\def\NomEvaluation{Examen 2 - Reprise}

% Ne rien mettre entre les accolades si l'enseignant (ou l'enseignante) est aussi le coordonnateur (ou la coordonnatrice)
\def\Coordonnatrice{}
\def\Coordonnateur{}

% Ne rien mettre entre les accolades s'il y a plus d'une section
\def\Enseignant{}
\def\Enseignante{Audrey-Anne Vallée}

% Ne rien mettre entre les accolades s'il s'agit d'un cours à section unique
% Mettre le nom de chacun des enseignants et enseignantes entre les accolades
% Une case à cocher sera générée pour chacune des sections
\def\SectionA{}
\def\SectionB{}
\def\SectionC{}
\def\SectionD{}


\def\Directives{
\begin{enumerate}
\item Matériel autorisé:
\begin{itemize}
\item Calculatrice scientifique {\bf autorisée par la FSG};
\item Le feuillet de formules est fourni avec l'examen. Celui-ci inclut les tables de lois. Veuillez ne rien écrire sur le feuillet de formules.
\end{itemize}
\item Assurez-vous que cette évaluation comporte \numquestions ~exercices répartis sur \numpages{}~pages, pour un total de \numpoints{}~points.\
\item \textbf{Veuillez ne pas dégrafer votre examen.}
\item Ajoutez vos initiales dans le bas de chaque page de l'examen.
\item Veuillez déposer votre carte d'identité \textbf{admise} avec photo sur votre bureau.\
\item {\bf Vous devez justifier chaque réponse, soit par une démarche, soit par une (des) phrase(s). Justifiez toutes les étapes de vos démarches.}
\item  Le verso des feuilles peut être utilisé comme brouillon, \textbf{mais ne sera pas corrigé}. Attention à ne rien dégrafer.
\end{enumerate}
}

\documentclass[Examen]{Evaluation}
\usepackage{multicol,graphicx}
\begin{document}

\begin{questions}  

\question 
Pour chaque énoncé qui suit, indiquez s'il est vrai ou faux.
\textbf{Justifiez} brièvement votre réponse. Il peut s'agir d'une justification mathématique ou d'une phrase. \\
\textit{Aucun point ne sera accordé pour une réponse sans justification ou une mauvaise justification.}

\vspace{0.3cm}
\begin{parts}
\part[4] La méthode de sélection de modèle pas-à-pas avec le critère BIC permet de trouver le modèle avec le plus petit BIC parmi tous les modèles possibles.

     \begin{oneparcheckboxes}
		\choice VRAI
		\correctchoice FAUX
\end{oneparcheckboxes}


\vspace{2.7 cm}

\part[4] Le coefficient de détermination ($R^2$) est un outil approprié pour faire la sélection de variables à inclure dans le modèle.

     \begin{oneparcheckboxes}
		\choice VRAI
		\correctchoice FAUX
\end{oneparcheckboxes}

\begin{solution}
Il augmente lorsque $p$ augmente, donc non.
\end{solution}

\vspace{2.7 cm}

\part[4]  Dans un modèle de régression linéaire multiple, la valeur-p du test F ajusté associé à une
variable explicative est toujours plus grande que celle du test t partiel de cette même variable.

     \begin{oneparcheckboxes}
		\choice VRAI
		\correctchoice FAUX
\end{oneparcheckboxes}

\begin{solution}
    Elles sont égales, les tests sont équivalents.
\end{solution}

\vspace{2.7 cm}

\part[4] Le levier d'une observation dépend en grande partie de sa valeur de la variable dépendante.

     \begin{oneparcheckboxes}
		\choice VRAI
		\correctchoice FAUX
\end{oneparcheckboxes}

\begin{solution}
$\frac{1}{1-R_j^2}$

$=1 si R_j^2 = 0$ et augmente si $R_j^2$ augmente.

Donc 1.
\end{solution}

\end{parts}


\newpage

\question
Pour chaque énoncé qui suit, justifiez brièvement votre réponse. Il peut s'agir d'une justification mathématique ou d'une phrase. \\

\begin{parts}
\part[4] Expliquez pourquoi il pourrait être préférable d'utiliser la distance de Cook plutôt que $DFBETA$.

\begin{solution}
    Ça fait augmenter la variance des estimateurs des coefficients de régression et des prévisions.
\end{solution}

\vspace{3 cm}

\part[4] Marie veut modéliser le temps requis pour que le fromage de sa fondue fonde en fonction de différentes variables explicatives. 
Elle dispose de 6 variables explicatives. Elle applique une sélection de modèle par exclusion basée sur le $R^2$ ajusté. À la fin de la première étape, combien de variables peut-elle avoir dans le modèle? Pour chaque réponse différente, expliquez comment on peut y arriver.
\begin{solution}
Elle n'identifie que les problèmes impliquant 2 variables explicatives, pas ceux impliquant plus de 2 variables explicatives.    
\end{solution}


\vspace{1.5 in }

\part[4] Soit le modèle de régression linéaire multiple pour lequel tous les postulats sont respectés. Calculez la covariance entre $\boldsymbol{Y}$ et $\boldsymbol{\varepsilon}$, c'est-à-dire $\mathrm{Cov}(\boldsymbol{y}, \boldsymbol{\varepsilon})$.

\begin{solution}
    $\mathrm{Cov}(y, \varepsilon) 
    = \mathrm{Cov}(X\beta+\varepsilon, \varepsilon)
    = \mathrm{Cov}(\varepsilon, \varepsilon)
    = = \mathrm{Var}(\varepsilon)
    = \sigma^2 \bf{1}$
\end{solution}

\vspace{3 cm}

\part[4] Que puis-je faire si j’ai un problème avec le postulat de linéarité et que je n’arrive pas à le corriger avec des transformations des variables? Donnez deux suggestions.

\begin{solution}
    Régression polynomiale, régression splines, arbres de régression
\end{solution}


\end{parts}


\newpage

\question\label{matrices}
On examine la réponse d'une variable $Y$ en fonction de deux variables explicatives $v$ et $w$, à partir d'un échantillon de taille $n=30$. Soit $\boldsymbol{X}=(\boldsymbol{1}, \boldsymbol{v}, \boldsymbol{w})$ une matrice $n\times 3$, où $\boldsymbol{1}$ est le vecteur de 1 associé à l'ordonnée à l'origine. Dans les résultats suivants, certaines valeurs de $\boldsymbol{X}^T\boldsymbol{X}$ ont été remplacées par les lettres $a, b, c$ et $d$. Pour les points (c) et (d), $\beta_v$ et $\beta_w$ sont les coefficients associés aux variables $\boldsymbol{v}$ et $\boldsymbol{w}$.
\bigskip

\begin{center}
    $\boldsymbol{X}^T\boldsymbol{X}=\left[
\begin{array}{ccc}
  a &20 & 0\\
 b & 20& 0\\
  c &d & 10\\
\end{array}\right];\qquad
(\boldsymbol{X}^T\boldsymbol{X})^{-1}=\left[
\begin{array}{ccc}
  0,1 &-0,1 & 0\\
 -0,1 & 0,15& 0\\
  0 &0 & 0,1\\
\end{array}\right];\quad
\boldsymbol{X}^T\boldsymbol{Y}=\begin{array}{c}
\end{array}\left[
\begin{array}{r}
  15 \\
  20 \\
 10\\
\end{array}\right];\quad
\boldsymbol{Y}^T\boldsymbol{Y}=59,5
$
\end{center}

\begin{parts}
    \part[3] Que valent $a, b$ et $c$? 

    \vspace{5.5cm}

    $a=$\boite{3}{1.5}
    $b=$\boite{3}{1.5}
    $c=$\boite{3}{1.5}
%    $d=$\boite{3}{1.5}

\part[3] Que valent $\overline{Y}$ et $\overline{v}$, les moyennes échantillonnales de $\boldsymbol{Y}$ et $\boldsymbol{v}$?

\vspace{5.5cm}
$\overline{Y}=$\boite{3}{1.5}
$\overline{v}=$\boite{3}{1.5}
%$\overline{w}=$\boite{3}{1.5}


\newpage

\part[3] En supposant que $\widehat{\sigma}^2=1$, que vaut $\widehat{\text{Var}}(\widehat{\beta}_v)$, l'estimateur de la variance de $\widehat{\beta}_v$?

\vspace{8cm}
$\widehat{\text{Var}}(\widehat{\beta}_1)=$\boite{3}{1.5}

\vspace{0.5cm}
\part[3] En supposant que $\widehat{\sigma}^2=1$, que vaut $\widehat{\text{Cov}}(\widehat{\beta}_v, \widehat{\beta}_w)$, l'estimateur de la covariance entre $\widehat{\beta}_v$ et $\widehat{\beta}_w$?

\vspace{8cm}
$\widehat{\text{Cov}}(\widehat{\beta}_v, \widehat{\beta}_w)=$\boite{3}{1.5}
\end{parts}


\newpage
\question\label{interactionEThypotheses}
Avec $i=1,\hdots,n$, nous voulons expliquer une variable réponse $Y_i$ à l'aide de $w_i$, une variable polytomique pouvant prendre les valeurs $\{A,B,C\}$, et la variable quantitative $x_i$. Le modèle de régression utilisé est 

\begin{equation}\label{eq:modeleInteraction}
    Y_i=\beta_0 + \beta_1 x_i + \beta_{A} a_{i} + \beta_B b_i + \beta_{IA}a_i x_i+ \beta_{IB}b_i x_i + \varepsilon_i,
\end{equation}
où 
\begin{align*}
    a_i =\left\{
			\begin{array}{ll}
			1  & \mbox{ si } w_i=A, \\
			0 & \mbox{ sinon, }
   \end{array}
			\right. & & b_i =\left\{
			\begin{array}{ll}
			1  & \mbox{ si } w_i=B, \\
			0 & \mbox{ sinon. }
   \end{array}
			\right. 
\end{align*}

Voici des énoncés basés sur ce modèle:

\begin{enumerate}
    \item Les trois catégories de la variable $w_i$ sont équivalentes pour prédire la variable $Y_i$, peu importe la valeur de la variable $x_i$;
    \item L'effet d'une augmentation de la variable $x_i$ est le même, peu importe la catégorie de la variable $w_i$;
    \item La variable $x_i$ n'a aucun effet sur la variable réponse;
    \item Il n'y a aucun lien linéaire entre la variable réponse et les variables explicatives (ni leur intéraction); 
    \item L'espérance de $Y_i$ est de $0$, lorsque $x_i=0$ et $a_1=1$.
\end{enumerate}

\vspace{0.25cm}
\begin{parts}
     \renewcommand*\half{,5}
\part[7\half] Pour chacun des énoncés ci-haut, cochez l'hypothèse qui correspond.
\begin{table}[h]
\begin{center}
\begin{tabular}{|c||ccccc|}
\hline
& \multicolumn{5}{c|}{Énoncés}\\ \hline
Hypothèses & \multicolumn{1}{l||}{~~~1~~~} & \multicolumn{1}{l||}{~~~2~~~} & \multicolumn{1}{l||}{~~~3~~~} & \multicolumn{1}{l||}{~~~4~~~} & \multicolumn{1}{l||}{~~~5~~~} \\ \hline
$\beta_A=\beta_B$ & \multicolumn{1}{l||}{}  & \multicolumn{1}{l||}{}  & \multicolumn{1}{l||}{}  & \multicolumn{1}{l||}{}  & \multicolumn{1}{l||}{}   \\ \hline
$\beta_A=\beta_B=\beta_{IA}=\beta_{IB}$ & \multicolumn{1}{l||}{}  & \multicolumn{1}{l||}{}  & \multicolumn{1}{l||}{}  & \multicolumn{1}{l||}{}  & \multicolumn{1}{l||}{}     \\ \hline
$\beta_{IA}=\beta_{IB}=0$ & \multicolumn{1}{l||}{}  & \multicolumn{1}{l||}{}  & \multicolumn{1}{l||}{}  & \multicolumn{1}{l||}{}  & \multicolumn{1}{l||}{}     \\ \hline
$\beta_1=\beta_{IA}=\beta_{IB}=0$ & \multicolumn{1}{l||}{}  & \multicolumn{1}{l||}{}  & \multicolumn{1}{l||}{}  & \multicolumn{1}{l||}{}  & \multicolumn{1}{l||}{}     \\ \hline
$\beta_A=\beta_B=\beta_{IA}=\beta_{IB}=0$ & \multicolumn{1}{l||}{}  & \multicolumn{1}{l||}{}  & \multicolumn{1}{l||}{}  & \multicolumn{1}{l||}{}  & \multicolumn{1}{l||}{}     \\ \hline
$\beta_1-\beta_{IA}-\beta_{IB}=0$ & \multicolumn{1}{l||}{}  & \multicolumn{1}{l||}{}  & \multicolumn{1}{l||}{}  & \multicolumn{1}{l||}{}  & \multicolumn{1}{l||}{}     \\ \hline
$\beta_{1}=0$ & \multicolumn{1}{l||}{}  & \multicolumn{1}{l||}{}  & \multicolumn{1}{l||}{}  & \multicolumn{1}{l||}{}  & \multicolumn{1}{l||}{}     \\ \hline
$\beta_1=\beta_A=\beta_B=\beta_{IA}=\beta_{IB}=0$ & \multicolumn{1}{l||}{}  & \multicolumn{1}{l||}{}  & \multicolumn{1}{l||}{}  & \multicolumn{1}{l||}{}  & \multicolumn{1}{l||}{}     \\ \hline
$\beta_0=\beta_1=\beta_A=\beta_B=\beta_{IA}=\beta_{IB}=0$ & \multicolumn{1}{l||}{}  & \multicolumn{1}{l||}{}  & \multicolumn{1}{l||}{}  & \multicolumn{1}{l||}{}  & \multicolumn{1}{l||}{}     \\ \hline
$\beta_0+\beta_A=0$ & \multicolumn{1}{l||}{}  & \multicolumn{1}{l||}{}  & \multicolumn{1}{l||}{}  & \multicolumn{1}{l||}{}  & \multicolumn{1}{l||}{}     \\ \hline
\end{tabular}
\end{center}
\end{table}
\newpage

\part[2\half] À partir du modèle (\ref{eq:modeleInteraction}), l'espérance de la variable réponse d'un individu augmente de combien d'unités si sa valeur de la variable quantitative $x$ vaut 3 unités de plus? Justifiez votre réponse.
% 5\beta_z + 5\beta_{ZT}t_i + 5\beta_{ZQ}q_i
\vspace{7cm}

\begin{solution}
    \begin{align*}
    \mathrm{E}(Y_i) &= \beta_0 + \beta_Z z_i + \beta_{T} t_{i} + \beta_Q q_i + \beta_{ZT}t_i z_i+ \beta_{ZQ}q_i z_i \cr 
    \mathrm{E}(Y_i|z_i+5) &= \beta_0 + \beta_Z (z_i+5) + \beta_{T} t_{i} + \beta_Q q_i + \beta_{ZT}t_i (z_i+5)+ \beta_{ZQ}q_i (z_i+5) \cr 
    diff &=  5\beta_Z  + 5\beta_{ZT}t_i + 5\beta_{ZQ}q_i 
\end{align*}
\end{solution}

\part[5] À partir du modèle (\ref{eq:modeleInteraction}), tracez un graphique de l'espérance de $Y$.
Notez que le graphique n'a pas besoin d'être parfait, ni à l'échelle.
Par contre, on doit y retrouver les éléments importants, les coefficients du modèle et le nom des axes. Supposez que tous les coefficients sont non nuls et de valeurs différentes.

\begin{solution}
Y ~ Z

3 droites

3 ordonnées

3 pentes    
\end{solution}
\end{parts}

\newpage

\question\label{demonstration}

Soit le mod\`ele de r\'egression lin\'eaire 
$$\boldsymbol{Y} = \boldsymbol{X} \boldsymbol{\beta} + {\boldsymbol{\varepsilon}},$$
o\`u $\mathrm{E}({\boldsymbol{\varepsilon}}) = {\bf 0}$ et $\mathrm{Var}({\boldsymbol{\varepsilon}})=\sigma^2 \bold{1}_n $, $\sigma>0$ et $\bold{1}_n$ est la matrice identit\'e de dimension $n\times n$. On considère l'estimateur  $$\widetilde{\boldsymbol{\beta}} = (\boldsymbol{X}^\top\boldsymbol{X} + \boldsymbol{\Omega} )^{-1} \boldsymbol{X}^\top\boldsymbol{Y},$$ 
où   $\boldsymbol{\Omega}$ est une matrice constante de dimension $p' \times p'$.

\begin{parts}
	\part[5]  Calculez l'espérance de $\widetilde{\boldsymbol{\beta}}$ et donnez son expression. Justifiez toutes les étapes de votre démarche.
 \vspace{2in}
	\part[1] À partir de l'espérance obtenue en (a), donnez la composition de la matrice $\boldsymbol{\Omega}$ pour laquelle  $\widetilde{\boldsymbol{\beta}}$ est un estimateur sans biais de  ${\boldsymbol{\beta}}$. Justifiez toutes les étapes de votre démarche.
 \vspace{2in}
    \part[6]  Calculez la variance de $\widetilde{\boldsymbol{\beta}}$ et donnez son expression. Justifiez toutes les étapes de votre démarche.
\end{parts}

%
\begin{solution}
	\begin{parts}
		\part $\mathrm{E}( \widetilde{\boldsymbol{\beta}} ) 
		= \mathrm{E}( (\boldsymbol{X}'\boldsymbol{X} + k \mathbb{1}_{p'} )^{-1} \boldsymbol{X}'\boldsymbol{Y})
		=  (\boldsymbol{X}'\boldsymbol{X} + k \mathbb{1}_{p'} )^{-1} \boldsymbol{X}'\mathrm{E}(\boldsymbol{Y})
		=  (\boldsymbol{X}'\boldsymbol{X} + k \mathbb{1}_{p'} )^{-1} \boldsymbol{X}'\boldsymbol{X}\boldsymbol{\beta}$
		\part Si $k=0$, $\mathrm{E}( \widetilde{\boldsymbol{\beta}} ) 
		= (\boldsymbol{X}'\boldsymbol{X} + k \mathbb{1}_{p'} )^{-1} \boldsymbol{X}'\boldsymbol{X}\boldsymbol{\beta}
		= (\boldsymbol{X}'\boldsymbol{X} + 0 \mathbb{1}_{p'} )^{-1} \boldsymbol{X}'\boldsymbol{X}\boldsymbol{\beta}
		= (\boldsymbol{X}'\boldsymbol{X} )^{-1} \boldsymbol{X}'\boldsymbol{X}\boldsymbol{\beta}
		=\boldsymbol{\beta}$
	\end{parts}
\end{solution}

\newpage




\question[8]\label{influence}

On veut étudier le lien linéaire entre 6 variables explicatives et une variable réponse. On dispose de la sortie statistique suivante:

\begin{verbatim}
Obs   Residual   RStudent       h_ii   CovRatio     DFFITS
1      -0.184      -0.075      0.149      1.580     -0.031
2       0.114       0.047      0.175      1.632      0.022
3       2.358       1.064      0.277      1.332      0.659
4      -3.377      -1.513      0.233      0.904     -0.833
5      -0.782      -2.986      0.987     10.190    -25.685

-------------------------------------DFBETAS---------------------------
Obs  Intercept      X1        X2        X3        X4        X5        X6
1      0.011     0.002    -0.007    -0.018     0.003    -0.014    -0.002
2      0.008    -0.005    -0.001     0.011    -0.007    -0.005    -0.006
3      0.332    -0.302     0.095    -0.301    -0.102    -0.322     0.009
4     -0.097     0.382    -0.267    -0.338     0.317    -0.165     0.065
5      0.581    -0.351    23.385     0.413    -0.642    -0.233    -6.733

\end{verbatim}

Il y a eu une erreur lors de la saisie de la valeur de \texttt{X2} de l'une des 5 premières observations (on a saisi 22.3 au lieu de 223). De quelle observation s'agit-il, justifiez votre choix? Est-ce que cette observation semble potentiellement influente pour certaines inférences? Justifiez vos réponses, des réponses sans justification recevront la note 0.

\newpage

\question\label{application}
On s'intéresse à la qualité de différents vins (mesurée par une note). On a collecté les variables suivantes (toutes continues) pour un échantillon aléatoire de vins:
\begin{center}
\begin{tabular}{lll}
     acidité fixe &  chlorures & pH \\
     acidité volatile &  anhydride sulfureux libre & sulfates \\
     acide citrique &   dioxyde de soufre total & alcool \\
     sucre résiduel & densité & 
\end{tabular}\end{center}
Vous faites l'hypothèses d'une relation linéaire entre la qualité du vin et les autres variables. Utilisez cette sortie \texttt{R}  pour répondre aux questions suivantes. {\bf Justifiez toutes vos réponses.}
\begin{verbatim}
Coefficients:
                       Estimate   Std. Error   t value   Pr(>|t|)    
(Intercept)             21.9700      21.1900     1.036     0.3002    
acidite.fixe             0.0250       0.0260     0.963     0.3357    
acidite.volatire        -1.0840       0.1211    -8.948    < 2e-16 
acide.citrice           -0.1826       0.1472    -1.240     0.2150    
sucre.residuel           0.0163       0.0150     1.089     0.2765    
chlorures               -1.8740       0.4193    -4.470   8.37e-06 
anhydrire.sulf.libre     0.0044       0.0022     2.009     0.0447 
dioxyde.soufre.total    -0.0033       0.0007    -4.480   8.00e-06 
densite                -17.8800      21.6300    -0.827     0.4086    
pH                      -0.4137       0.1916    -2.159     0.0310 
sulfates                 0.9163       0.1143     8.014   2.13e-15 
alcool                   0.2762       0.2648    10.429    < 2e-16 

Residual standard error:   ?   on  1587  degrees of freedom
Multiple R-squared:    ? ,	Adjusted R-squared:   0.3561 
F-statistic:  81.35  on  11  and  1587  DF,  p-value:  < 2.2e-16    
\end{verbatim}

\begin{parts}
    \part[2] Combien de vins font partie de l'échantillon?
    \vspace{1in}
    \part[4] On pense qu'un vin avec une quantité d'anhydride sulfureux libre (\texttt{anhydrire.sulf.libre}) supérieure à celle d'un autre vin (toutes autres variables explicatives étant égales) a en moyenne une qualité supérieure. Au seuil 5\%, est-ce que les résultats du modèle confirment cette hypothèse?
    \vspace{2in}
    \part[4] Calculez le pourcentage de la variabilité de la qualité du vin qui est expliquée par le modèle.
    \vspace{3in}
    \part[7] Lorsque vous magasinez votre vin, vous aimeriez choisir des vins en fonction de leur qualité. Vous aimeriez construire le meilleur modèle de régression possible afin de prédire avec justesse la valeur de la qualité en fonction des variables exogènes disponibles. Après quelques manipulations, vous obtenez les résultats suivants:
	\begin{center}\begin{tabular}{|c|c|c|c|c|c|}\hline
		Mod\`ele & $R^2$ & $R_{a}^2$ & AIC & BIC & $PRESS$\\ \hline\hline
		Modèle 1 & 0.3606 & 0.3561 & 3164.277 & 3234.179 & 679.2157 \\\hline
            Modèle 2 & 0.3603 & 0.3562 & 3162.965 & 3227.490 & 678.3656 \\\hline
            Modèle 3 & 0.3602 & 0.3565 & 3161.224 & 3220.372 & 677.3349 \\\hline
            Modèle 4 & 0.3599 & 0.3567 & 3158.839 & 3207.610 & 676.3202 \\\hline
            Modèle 5 & 0.3595 & 0.3566 & 3159.977 & 3213.371 & 675.8162 \\\hline
            Modèle 6 & 0.3594 & 0.3527 & 3173.829 & 3216.847 & 681.8370 \\\hline
	\end{tabular}\end{center}
    Sur la base de ce tableau, quel modèle est préférable pour la mise en situation de cette sous-question? Justifiez votre réponse, un choix de modèle non justifié ne recevra aucun points. 
    \newpage
    \part[4] On ajuste un modèle avec trois variables explicatives uniquement. Voici une sortie obtenue avec ce modèle:
    \begin{verbatim}
  Eigenvalue  Condition Index   intercept        x1        x2        x3
1     3.4486           1.0000      0.0110    0.0152    0.0123    0.0133
2     0.3456           3.1589      0.0359    0.2965    0.1019    0.1311
3     0.1122           5.5442      0.2722    0.2247    0.3764    0.6967
4     0.0936           6.0699      0.6809    0.4636    0.5095    0.1589
    \end{verbatim}
    Interprétez cette sortie. Est-ce qu'un problème est détecté? Si oui, expliquez quel est le problème, quelles variables sont impliquées et ce que vous feriez par la suite. Si non, expliquez pourquoi il ne semble pas y avoir de problème. Dans vos explications, rapportez des nombres de la sortie pour justifier vos conclusions.
\end{parts}


\end{questions}

\end{document}
\newpage
    \part[6] Voici de l'information sur l'observation avec le plus grand résidu studentisé en valeur absolue.
    
    \begin{verbatim}
    h     RStudent     DFFITS    CovRatio     Cook
0.011       -4.194     -0.438       0.892    0.016
    \end{verbatim}
    Répondez aux questions suivantes. Des réponses non justifiées ne recevront aucun points.
    
    \begin{subparts}
%        \item Est-ce que cette observation a de l'influence sur l'estimation des coefficients de régression? \vspace{2in}
        \subpart Est-ce que cette observation fait augmenter ou diminuer la variance du vecteur de coefficients de régression estimés?\vspace{2.5in}
        \subpart Est-ce que la présence de cette observation fait augmenter ou diminuer la valeur de sa propre prédiction? Quantifiez cette augmentation ou diminution.\vspace{2in}
    \end{subparts}
    
